%%%%%%%%%%%%%%%%%%%%%%%%%%%%%%%%
% Chapter 4. Results and Discussion
% 예시: 표(booktabs + multirow), subfigure 환경을 이용한 부분 그림
%%%%%%%%%%%%%%%%%%%%%%%%%%%%%%%%
\chapter{Results and Discussion} \label{chapter4}

%==============================%
% 4.1 Results
%==============================%
\section{Results} \label{chap4:sec1}

\subsection{Quantitative results} \label{chap4:sec1:subsec1}

\Cref{tab:results} compares the performance of the models. Tables with grouped rows and columns can be built with \texttt{booktabs}, \texttt{multirow}, and \textbackslash\texttt{cmidrule}.

%----------------------------------------------------------
% [표 예시 2] tabular + booktabs + multirow
%  - \multirow{2}{*}{...} : 세로 2칸 병합
%  - \multicolumn{2}{c}{...} : 가로 2칸 병합
%  - \cmidrule(lr){3-4} : 3~4열에만 가로선
%----------------------------------------------------------
\begin{table}[!htbp]
\centering
\small
\caption{Performance comparison (example of a \texttt{booktabs} table).}
\label{tab:results}
\setlength{\tabcolsep}{6pt}
\begin{tabular}{cc cc cc}
\toprule
\multirow{2}{*}{\textbf{Model}} & \multirow{2}{*}{\textbf{Input}} &
\multicolumn{2}{c}{\textbf{Site A}} &
\multicolumn{2}{c}{\textbf{Site B}} \\
\cmidrule(lr){3-4} \cmidrule(lr){5-6}
& & Metric 1 & Metric 2 & Metric 1 & Metric 2 \\
\midrule
\multirow{2}{*}{Model 1}
& Baseline      & 0.00 & 0.00 & 0.00 & 0.00 \\
& \textbf{Proposed} & 0.00 & 0.00 & \textbf{0.00} & \textbf{0.00} \\
\addlinespace
\multirow{2}{*}{Model 2}
& Baseline      & 0.00 & 0.00 & 0.00 & 0.00 \\
& \textbf{Proposed} & 0.00 & 0.00 & \textbf{0.00} & \textbf{0.00} \\
\bottomrule
\end{tabular}
\end{table}


\subsection{Qualitative results} \label{chap4:sec1:subsec2}

% ------------------------------------------------------------------
% [예시 4] 그림 파일이 패널별로 따로 있을 때: subfigure 환경
% 각 패널 아래에 (a), (b) 캡션이 자동으로 붙습니다.
% ------------------------------------------------------------------
\begin{figure}[!htbp]
    \centering
    \begin{subfigure}[b]{0.48\textwidth}
        \centering
        \includegraphics[width=\textwidth]{example-image-b}
        \caption{Baseline}
        \label{fig:map_baseline}
    \end{subfigure}
    \hfill
    \begin{subfigure}[b]{0.48\textwidth}
        \centering
        \includegraphics[width=\textwidth]{example-image-c}
        \caption{Proposed}
        \label{fig:map_proposed}
    \end{subfigure}
    \caption{Visual comparison of results: (a) baseline and (b) proposed method.}
    \label{fig:maps}
\end{figure}

\Cref{fig:maps} compares the results visually. The proposed method (\Cref{fig:map_proposed}) shows ... compared with the baseline (\Cref{fig:map_baseline}).

%==============================%
% 4.2 Discussion
%==============================%
\section{Discussion} \label{chap4:sec2}

Interpret the results, compare them with previous studies, and discuss their implications.

\subsection{Limitations and future work} \label{chap4:sec2:subsec1}

Describe the limitations of this study and directions for future research.
